% ============================================================================
%  RAPPORT DE BENCHMARK DE CALIBRAGE — MISSION D'AUDIT GSE
%  Genere automatiquement par audit_mission/benchmark/report.py
%  Compilation : pdflatex (x2) ou latexmk -pdf   (depuis audit_mission/)
% ============================================================================
\documentclass[11pt,a4paper]{article}
\usepackage[utf8]{inputenc}
\usepackage[T1]{fontenc}
\usepackage{helvet}
\renewcommand{\familydefault}{\sfdefault}
\frenchspacing
\providecommand{\og}{\guillemotleft~}
\providecommand{\fg}{~\guillemotright}
\usepackage[a4paper,margin=2cm,headheight=30pt,headsep=18pt,footskip=28pt]{geometry}
\usepackage{setspace}\setstretch{1.08}
\setlength{\parindent}{0pt}\setlength{\parskip}{5pt}
\usepackage{xcolor,colortbl}
\definecolor{nexred}{HTML}{C8102E}
\definecolor{nexnavy}{HTML}{14213D}
\definecolor{nexgrey}{HTML}{808080}
\definecolor{nexlight}{HTML}{F2F4F8}
\definecolor{nexgreen}{HTML}{1E7B4F}
\definecolor{tabborder}{HTML}{BFC5D2}
\usepackage{booktabs,tabularx,array,float,graphicx}
\usepackage{amsmath,amssymb}
\usepackage{caption}
\usepackage{longtable,multirow}
\usepackage{ragged2e,enumitem}
\setlist[itemize]{label=--, leftmargin=1.2em, itemsep=2pt, topsep=3pt}
\usepackage{fancyhdr,lastpage}
\usepackage{titlesec}
\usepackage[most]{tcolorbox}
\usepackage[hidelinks]{hyperref}
\titleformat{\section}{\LARGE\bfseries\color{nexnavy}}{\thesection.}{0.6em}{}
  [{\color{nexred}\titlerule[1.4pt]}]
\titleformat{\subsection}{\large\bfseries\color{nexnavy}}{\thesubsection}{0.6em}{}
\titleformat{\subsubsection}{\normalsize\bfseries\color{nexred}}{\thesubsubsection}{0.6em}{}
\titlespacing*{\section}{0pt}{18pt}{10pt}
\titlespacing*{\subsection}{0pt}{14pt}{6pt}
\fancypagestyle{nexialog}{%
  \fancyhf{}
  \fancyhead[L]{\small\textbf{\textcolor{nexred}{NEXIALOG}\,\textcolor{nexnavy}{CONSULTING}}}
  \fancyhead[R]{\small\textcolor{nexgrey}{Benchmark de calibrage GSE \,|\, Confidentiel}}
  \fancyfoot[L]{\footnotesize\textcolor{nexgrey}{\copyright~Nexialog Consulting}}
  \fancyfoot[R]{\footnotesize Page \thepage~/~\pageref{LastPage}}
  \renewcommand{\headrulewidth}{1.2pt}
  \renewcommand{\headrule}{\color{nexred}\hrule width\headwidth height\headrulewidth}
  \renewcommand{\footrulewidth}{0.4pt}}
\pagestyle{nexialog}
\newcolumntype{P}[1]{>{\RaggedRight\arraybackslash}p{#1}}
\arrayrulecolor{tabborder}
\renewcommand{\arraystretch}{1.35}
\newcommand{\thead}[1]{\textcolor{white}{\textbf{\small #1}}}
\newcommand{\tabcaption}[1]{\par\vspace{2pt}{\centering\footnotesize\textcolor{nexgrey}{\itshape #1}\par}\vspace{8pt}}
\newtcolorbox{findingbox}[1]{breakable, colback=nexlight, colframe=nexred,
  boxrule=0.8pt, left=8pt, right=8pt, top=6pt, bottom=6pt, arc=2pt,
  fonttitle=\bfseries\color{white}, coltitle=white, title={#1}}
\begin{document}


\begin{titlepage}
  \thispagestyle{empty}
  \vspace*{2.0cm}
  {\large\bfseries\color{nexred} RAPPORT DE BENCHMARK DE CALIBRAGE\par}
  \vspace{0.4cm}
  {\setstretch{1.15}\Huge\bfseries\color{nexnavy}
   Confrontation ind\'ependante du calibrage du G\'en\'erateur de Sc\'enarios
   \'Economiques : notre outil vs param\`etres en production\par}
  \vspace{0.6cm}
  {\large\color{nexgrey} Mission d'audit GSE --- P\^ole Risques Financiers \& Actuariat\par}
  \vspace{0.4cm}
  {\normalsize\color{nexnavy}\itshape Exercice d'allocation 2026 --- param\`etres GSE arr\^et\'es au 31/12/2025\par}
  \vspace{1.0cm}{\color{nexred}\hrule height 1.6pt}\vspace{0.5cm}
  Fen\^etre et pas : \textbf{variables selon le facteur} (d\'etaill\'es en t\^ete de chaque section).\par
  Param\`etres en production confront\'es : millésime 2026 (20 ans).\par
  \vfill
\end{titlepage}
\tableofcontents
\clearpage

\section{Objet et p\'erim\`etre}
Ce rapport restitue, facteur de risque par facteur de risque, la confrontation
entre (i) le calibrage obtenu avec \textbf{notre propre outil}
(\texttt{gse\_engine.ahlgrim}) sur les donn\'ees de calibrage fournies, et (ii)
les \textbf{param\`etres d\'ej\`a calibr\'es} transmis par le client (classeur
\texttt{Parametres\_models.xlsx}, mill\'esime 2026). Pour
chaque facteur sont pr\'esent\'es : un tableau de comparaison des param\`etres,
un tableau de comparaison des m\'etriques de validation, et les figures de
diagnostic (QQ-plot et autocorr\'elogramme des r\'esidus standardis\'es).

\subsection{D\'emarche}
\begin{itemize}
  \item \textbf{Calibrage (notre outil)} : chaque mod\`ele du moteur est
        calibr\'e sur les donn\'ees fournies, sur la fen\^etre et le pas \textbf{propres \`a chaque facteur} (cf. sections).
        En sortent les param\`etres et les r\'esidus.
  \item \textbf{Param\`etres en production} : les param\`etres du classeur sont
        inject\'es dans le m\^eme moteur ; les r\'esidus sont recalcul\'es sur les
        m\^emes donn\'ees, puis les m\'etriques et figures.
  \item \textbf{M\'etriques} : $R^2$ et $R^2$ ajust\'e, Shapiro-Wilk et
        Jarque-Bera (normalit\'e), Ljung-Box et ACF (autocorr\'elation). Les
        r\'esidus standardis\'es sont $(\varepsilon-\bar\varepsilon)/s_\varepsilon$ ;
        les tests de normalit\'e et d'autocorr\'elation \'etant invariants
        d'\'echelle, les conclusions ne d\'ependent pas de l'\'echelle du $\sigma$ (donc
        de la coquille de facteur 10, constat U-1).
\end{itemize}

\subsection{Correspondance mod\`eles / donn\'ees}
\begin{itemize}
  \item Inflation, taux r\'eels : Vasicek \`a deux facteurs (donn\'ees \og rate \fg, niveaux) ;
  \item Cr\'edit IG : CIR (\og rate \fg) ; dette priv\'ee : Black-Karasinski / log-OU
        (s\'erie d\'ej\`a en $100\log$, OU direct) ;
  \item Actions cot\'ees (Euro, Monde, \'Emergent) : Hardy RSLN-2 (donn\'ees \og brutes \fg, log-rendements) ;
  \item Actions non cot\'ees, immobilier, infrastructure : Black-Scholes apr\`es
        \emph{unsmoothing} AR(1) (donn\'ees \og brutes \fg).
\end{itemize}

\subsection{Constats transverses de calibrage}
\begin{findingbox}{Constat U-1 --- Coquille de facteur 10 sur les volatilit\'es (identifi\'ee puis corrig\'ee \`a la source)}
\textbf{Anomalie identifi\'ee, puis corrig\'ee par le client.} Les volatilit\'es
du classeur pour les mod\`eles sur log-rendements (Hardy, Black-Scholes) et pour
l'OU sur log-spread (dette priv\'ee) \'etaient initialement \textbf{amput\'ees
d'un facteur 10} par rapport \`a l'unit\'e du moteur (\'etat $z=100\log$). La
preuve \'etait exacte sur le Private Equity : notre s\'erie deslissée donne une
volatilit\'e post-unsmoothing de $29{,}70$ alors que le classeur indiquait
$2{,}970$. \textbf{La coquille a \'et\'e corrig\'ee directement dans le classeur
source} (\texttt{Parametres\_models.xlsx}) ; le pr\'esent benchmark utilise le
classeur corrig\'e et \textbf{le code n'applique aucune correction}
($\sigma$ pris tel quel). \textbf{V\'erification a posteriori} : apr\`es
correction, le recalibrage co\"incide avec la production \`a $\sim$1--2\,\%
pr\`es (ex. actions Euro $\sigma_1$ : $21{,}48$ vs $21{,}76$ ; infrastructure
$\sigma$ : $13{,}38$ vs $13{,}50$).

\textbf{Périm\`etre.} La coquille ne concernait que les volatilit\'es sur
log-rendements / log-spread. La volatilit\'e du CIR cr\'edit ($0{,}517$)
co\"incide d\'ej\`a avec notre recalibrage ($0{,}534$) : exprim\'ee sur
l'\'echelle de diffusion du spread ($\sigma\sqrt{s}$), elle n'\'etait pas
concern\'ee (un facteur 10 y donnerait une dispersion aberrante). Les
volatilit\'es V2F (taux, inflation) rel\`event d'une diff\'erence
m\'ethodologique distincte (constat M-1).
\end{findingbox}
\begin{findingbox}{Constat M-1 --- M\'ethodologie V2F (taux / inflation)}
Notre outil calibre le V2F par \textbf{maximum de vraisemblance} ; la production
calibre par \textbf{cibles distributionnelles}. Les vitesses de retour
($\kappa$) et la moyenne long terme sont proches ; les volatilit\'es
diff\`erent sensiblement par construction (volatilit\'e d'innovation MV vs
volatilit\'e stationnaire cible). La comparaison des r\'esidus reste valide
(la pr\'ediction conditionnelle ne d\'epend que des $\kappa$ et de $\mu_l$).
\end{findingbox}
\begin{findingbox}{Constat D-1 --- Dérive immobilier (rendement total vs prix)}
La d\'erive $\mu$ immobilier en production int\`egre le r\'einvestissement des
loyers (rendement total), alors que la s\'erie fournie est un indice de prix.
L'\'ecart de $\mu$ est donc attendu et ne traduit pas une erreur de calibrage.
\end{findingbox}

\clearpage
\section{Synth\`ese}
Le tableau ci-dessous r\'esume, pour chaque facteur, le nombre d'observations,
le $R^2$ (notre outil / production), la p-value de Shapiro-Wilk et le verdict
d'ind\'ependance (Ljung-Box, seuil 5\,\%).
\input{benchmark_output/tables/synthese.tex}
\tabcaption{Synth\`ese du benchmark de calibrage (notre outil vs production 2026).}
\clearpage

\section{Inflation (Vasicek 2 facteurs)}
\subsection{Mod\`ele et donn\'ees}
\textbf{Mod\`ele :} Vasicek à deux facteurs (V2F).
\begin{equation*}
dq^s_t=\kappa_1(q^l_t-q^s_t)\,dt+\sigma_1\,dW^1_t,\qquad dq^l_t=\kappa_2(\mu_l-q^l_t)\,dt+\sigma_2\,dW^2_t
\end{equation*}
\textbf{Fen\^etre / pas :} 31/12/2005 -- 31/12/2025 \`a pas mensuel ($\delta=1/12$).
\textbf{Donn\'ees :} Inflation point mort (BEIR) 2 ans (facteur court) et 10 ans (facteur long latent), à pas mensuel.
\textbf{Note de calibrage :} Notre outil calibre le V2F par maximum de vraisemblance exact (discrétisation OU bivariée). Le dispositif en production calibre par cibles distributionnelles (volatilité des variations, corrélation inter-maturités, écart-type des taux) : les volatilités diffèrent donc par construction méthodologique. La corrélation $\rho$ des browniens est fixée à 0 en production (remédiation), retenue ici pour le recalcul des résidus.
\subsection{Comparaison des param\`etres}
\input{benchmark_output/tables/inflation_params.tex}
\tabcaption{Param\`etres calibr\'es : notre outil vs production 2026 --- Inflation (Vasicek 2 facteurs).}
\subsection{M\'etriques de validation}
\input{benchmark_output/tables/inflation_metrics.tex}
\tabcaption{M\'etriques de validation des r\'esidus --- Inflation (Vasicek 2 facteurs).}
\subsection{Figures de diagnostic}
\begin{figure}[H]\centering
\includegraphics[width=\linewidth]{benchmark_output/figures/inflation_qq.pdf}
\caption*{\footnotesize\itshape QQ-plot des r\'esidus standardis\'es vs loi normale --- Inflation (Vasicek 2 facteurs) (gauche : notre outil ; droite : production).}
\end{figure}
\begin{figure}[H]\centering
\includegraphics[width=\linewidth]{benchmark_output/figures/inflation_acf.pdf}
\caption*{\footnotesize\itshape Autocorr\'elogramme (ACF) des r\'esidus (retards successifs, bande IC 95\,\%) --- Inflation (Vasicek 2 facteurs).}
\end{figure}
\clearpage
\section{Taux reels (Vasicek 2 facteurs)}
\subsection{Mod\`ele et donn\'ees}
\textbf{Mod\`ele :} Vasicek à deux facteurs (V2F).
\begin{equation*}
dr_t=\kappa_1(l_t-r_t)\,dt+\sigma_1\,dW^1_t,\qquad dl_t=\kappa_2(\mu-l_t)\,dt+\sigma_2\,dW^2_t
\end{equation*}
\textbf{Fen\^etre / pas :} 31/12/2005 -- 31/12/2025 \`a pas mensuel ($\delta=1/12$).
\textbf{Donn\'ees :} Taux réels 2 ans et 10 ans (reconstruits par Fisher), pas mensuel.
\textbf{Note de calibrage :} Mêmes remarques méthodologiques que pour l'inflation (MV vs cibles distributionnelles, $\rho=0$).
\subsection{Comparaison des param\`etres}
\input{benchmark_output/tables/real_rate_params.tex}
\tabcaption{Param\`etres calibr\'es : notre outil vs production 2026 --- Taux reels (Vasicek 2 facteurs).}
\subsection{M\'etriques de validation}
\input{benchmark_output/tables/real_rate_metrics.tex}
\tabcaption{M\'etriques de validation des r\'esidus --- Taux reels (Vasicek 2 facteurs).}
\subsection{Figures de diagnostic}
\begin{figure}[H]\centering
\includegraphics[width=\linewidth]{benchmark_output/figures/real_rate_qq.pdf}
\caption*{\footnotesize\itshape QQ-plot des r\'esidus standardis\'es vs loi normale --- Taux reels (Vasicek 2 facteurs) (gauche : notre outil ; droite : production).}
\end{figure}
\begin{figure}[H]\centering
\includegraphics[width=\linewidth]{benchmark_output/figures/real_rate_acf.pdf}
\caption*{\footnotesize\itshape Autocorr\'elogramme (ACF) des r\'esidus (retards successifs, bande IC 95\,\%) --- Taux reels (Vasicek 2 facteurs).}
\end{figure}
\clearpage
\section{Credit - spread IG (CIR)}
\subsection{Mod\`ele et donn\'ees}
\textbf{Mod\`ele :} Cox-Ingersoll-Ross (CIR), schéma d'Alfonsi.
\begin{equation*}
ds_t=\kappa(\theta-s_t)\,dt+\sigma\sqrt{s_t}\,dW_t
\end{equation*}
\textbf{Fen\^etre / pas :} 31/12/2005 -- 31/12/2025 \`a pas mensuel ($\delta=1/12$).
\textbf{Donn\'ees :} Spread Investment Grade bucket 5--7 ans, pas mensuel.
\textbf{Note de calibrage :} Calibrage par maximum de vraisemblance sur la pseudo-vraisemblance gaussienne d'Alfonsi. Les résidus sont standardisés par la volatilité conditionnelle $\sigma\sqrt{s_{t-1}\delta}$. La condition de Feller $2\kappa\theta\geq\sigma^2$ est contrôlée.
\subsection{Comparaison des param\`etres}
\input{benchmark_output/tables/credit_params.tex}
\tabcaption{Param\`etres calibr\'es : notre outil vs production 2026 --- Credit - spread IG (CIR).}
\subsection{M\'etriques de validation}
\input{benchmark_output/tables/credit_metrics.tex}
\tabcaption{M\'etriques de validation des r\'esidus --- Credit - spread IG (CIR).}
\subsection{Figures de diagnostic}
\begin{figure}[H]\centering
\includegraphics[width=\linewidth]{benchmark_output/figures/credit_qq.pdf}
\caption*{\footnotesize\itshape QQ-plot des r\'esidus standardis\'es vs loi normale --- Credit - spread IG (CIR) (gauche : notre outil ; droite : production).}
\end{figure}
\begin{figure}[H]\centering
\includegraphics[width=\linewidth]{benchmark_output/figures/credit_acf.pdf}
\caption*{\footnotesize\itshape Autocorr\'elogramme (ACF) des r\'esidus (retards successifs, bande IC 95\,\%) --- Credit - spread IG (CIR).}
\end{figure}
\clearpage
\section{Dette privee (Black-Karasinski / log-OU)}
\subsection{Mod\`ele et donn\'ees}
\textbf{Mod\`ele :} Black-Karasinski (log-OU).
\begin{equation*}
d\ln s_t=k\,(\theta-\ln s_t)\,dt+\sigma\,dW_t
\end{equation*}
\textbf{Fen\^etre / pas :} 31/12/2005 -- 31/12/2025 \`a pas mensuel ($\delta=1/12$).
\textbf{Donn\'ees :} Spread de dette privée déjà transformé en $100\log(s_t)$ : l'OU est appliqué directement sur la série fournie.
\textbf{Note de calibrage :} L'initialisation de la vraisemblance OU est guidée par les données (moyenne $\approx-305$) afin d'assurer la convergence. Les volatilités du classeur (corrigé de la coquille de facteur 10, cf. constat U-1) sont en unité moteur.
\subsection{Comparaison des param\`etres}
\input{benchmark_output/tables/dette_privee_params.tex}
\tabcaption{Param\`etres calibr\'es : notre outil vs production 2026 --- Dette privee (Black-Karasinski / log-OU).}
\subsection{M\'etriques de validation}
\input{benchmark_output/tables/dette_privee_metrics.tex}
\tabcaption{M\'etriques de validation des r\'esidus --- Dette privee (Black-Karasinski / log-OU).}
\subsection{Figures de diagnostic}
\begin{figure}[H]\centering
\includegraphics[width=\linewidth]{benchmark_output/figures/dette_privee_qq.pdf}
\caption*{\footnotesize\itshape QQ-plot des r\'esidus standardis\'es vs loi normale --- Dette privee (Black-Karasinski / log-OU) (gauche : notre outil ; droite : production).}
\end{figure}
\begin{figure}[H]\centering
\includegraphics[width=\linewidth]{benchmark_output/figures/dette_privee_acf.pdf}
\caption*{\footnotesize\itshape Autocorr\'elogramme (ACF) des r\'esidus (retards successifs, bande IC 95\,\%) --- Dette privee (Black-Karasinski / log-OU).}
\end{figure}
\clearpage
\section{Action Euro (Hardy RSLN-2)}
\subsection{Mod\`ele et donn\'ees}
\textbf{Mod\`ele :} Hardy (RSLN-2, changement de régime).
\begin{equation*}
z_t\mid\rho_t=i\ \sim\ \mathcal N(\delta\mu_i,\ \delta\sigma_i^2),\quad i\in\{1,2\}
\end{equation*}
\textbf{Fen\^etre / pas :} 31/12/2005 -- 31/12/2025 \`a pas mensuel ($\delta=1/12$).
\textbf{Donn\'ees :} Indice actions zone euro, log-rendements mensuels.
\textbf{Note de calibrage :} Calibrage EM (Baum-Welch). Le $R^2$ est structurellement proche de 0 (la moyenne conditionnelle de régime explique peu la variance des rendements) : les tests pertinents sont la normalité conditionnelle (SW/QQ) et l'absence d'autocorrélation. Les volatilités de régime du classeur (corrigé de la coquille de facteur 10, cf. constat U-1) sont en unité moteur.
\subsection{Comparaison des param\`etres}
\input{benchmark_output/tables/action_eur_params.tex}
\tabcaption{Param\`etres calibr\'es : notre outil vs production 2026 --- Action Euro (Hardy RSLN-2).}
\subsection{M\'etriques de validation}
\input{benchmark_output/tables/action_eur_metrics.tex}
\tabcaption{M\'etriques de validation des r\'esidus --- Action Euro (Hardy RSLN-2).}
\subsection{Figures de diagnostic}
\begin{figure}[H]\centering
\includegraphics[width=\linewidth]{benchmark_output/figures/action_eur_qq.pdf}
\caption*{\footnotesize\itshape QQ-plot des r\'esidus standardis\'es vs loi normale --- Action Euro (Hardy RSLN-2) (gauche : notre outil ; droite : production).}
\end{figure}
\begin{figure}[H]\centering
\includegraphics[width=\linewidth]{benchmark_output/figures/action_eur_acf.pdf}
\caption*{\footnotesize\itshape Autocorr\'elogramme (ACF) des r\'esidus (retards successifs, bande IC 95\,\%) --- Action Euro (Hardy RSLN-2).}
\end{figure}
\clearpage
\section{Action Monde ex-Europe (Hardy RSLN-2)}
\subsection{Mod\`ele et donn\'ees}
\textbf{Mod\`ele :} Hardy (RSLN-2, changement de régime).
\begin{equation*}
z_t\mid\rho_t=i\ \sim\ \mathcal N(\delta\mu_i,\ \delta\sigma_i^2)
\end{equation*}
\textbf{Fen\^etre / pas :} 31/12/2005 -- 31/12/2025 \`a pas mensuel ($\delta=1/12$).
\textbf{Donn\'ees :} Indice actions Monde ex-Europe, log-rendements mensuels.
\textbf{Note de calibrage :} Mêmes remarques que pour les actions Euro (constat U-1 sur les volatilités de régime).
\subsection{Comparaison des param\`etres}
\input{benchmark_output/tables/action_monde_params.tex}
\tabcaption{Param\`etres calibr\'es : notre outil vs production 2026 --- Action Monde ex-Europe (Hardy RSLN-2).}
\subsection{M\'etriques de validation}
\input{benchmark_output/tables/action_monde_metrics.tex}
\tabcaption{M\'etriques de validation des r\'esidus --- Action Monde ex-Europe (Hardy RSLN-2).}
\subsection{Figures de diagnostic}
\begin{figure}[H]\centering
\includegraphics[width=\linewidth]{benchmark_output/figures/action_monde_qq.pdf}
\caption*{\footnotesize\itshape QQ-plot des r\'esidus standardis\'es vs loi normale --- Action Monde ex-Europe (Hardy RSLN-2) (gauche : notre outil ; droite : production).}
\end{figure}
\begin{figure}[H]\centering
\includegraphics[width=\linewidth]{benchmark_output/figures/action_monde_acf.pdf}
\caption*{\footnotesize\itshape Autocorr\'elogramme (ACF) des r\'esidus (retards successifs, bande IC 95\,\%) --- Action Monde ex-Europe (Hardy RSLN-2).}
\end{figure}
\clearpage
\section{Action Emergent (Hardy RSLN-2)}
\subsection{Mod\`ele et donn\'ees}
\textbf{Mod\`ele :} Hardy (RSLN-2, changement de régime).
\begin{equation*}
z_t\mid\rho_t=i\ \sim\ \mathcal N(\delta\mu_i,\ \delta\sigma_i^2)
\end{equation*}
\textbf{Fen\^etre / pas :} 31/12/2005 -- 31/12/2025 \`a pas mensuel ($\delta=1/12$).
\textbf{Donn\'ees :} Indice actions Émergents, log-rendements mensuels.
\textbf{Note de calibrage :} Mêmes remarques que pour les actions Euro (constat U-1).
\subsection{Comparaison des param\`etres}
\input{benchmark_output/tables/action_em_params.tex}
\tabcaption{Param\`etres calibr\'es : notre outil vs production 2026 --- Action Emergent (Hardy RSLN-2).}
\subsection{M\'etriques de validation}
\input{benchmark_output/tables/action_em_metrics.tex}
\tabcaption{M\'etriques de validation des r\'esidus --- Action Emergent (Hardy RSLN-2).}
\subsection{Figures de diagnostic}
\begin{figure}[H]\centering
\includegraphics[width=\linewidth]{benchmark_output/figures/action_em_qq.pdf}
\caption*{\footnotesize\itshape QQ-plot des r\'esidus standardis\'es vs loi normale --- Action Emergent (Hardy RSLN-2) (gauche : notre outil ; droite : production).}
\end{figure}
\begin{figure}[H]\centering
\includegraphics[width=\linewidth]{benchmark_output/figures/action_em_acf.pdf}
\caption*{\footnotesize\itshape Autocorr\'elogramme (ACF) des r\'esidus (retards successifs, bande IC 95\,\%) --- Action Emergent (Hardy RSLN-2).}
\end{figure}
\clearpage
\section{Action non cotee - PE (Black-Scholes apres unsmoothing)}
\subsection{Mod\`ele et donn\'ees}
\textbf{Mod\`ele :} Black-Scholes après unsmoothing AR(1).
\begin{equation*}
\frac{dV_t}{V_t}=\mu\,dt+\sigma\,dW_t\quad\text{(série deslissée)}
\end{equation*}
\textbf{Fen\^etre / pas :} 31/12/2005 -- 31/12/2025 \`a pas mensuel ($\delta=1/12$).
\textbf{Donn\'ees :} Indice LPX50 (Private Equity coté), log-rendements mensuels deslissés.
\textbf{Note de calibrage :} Le deslissage AR(1) (régression $r_t=a+b\,r_{t-1}+e_t$, série $r^{uns}=(r_t-b\,r_{t-1})/(1-b)$) reproduit exactement la volatilité production ($29{,}70$), portée au classeur corrigé (cf. constat U-1). L'ACF(1) post-unsmoothing doit être proche de 0.
\subsection{Comparaison des param\`etres}
\input{benchmark_output/tables/action_noncotee_params.tex}
\tabcaption{Param\`etres calibr\'es : notre outil vs production 2026 --- Action non cotee - PE (Black-Scholes apres unsmoothing).}
\subsection{M\'etriques de validation}
\input{benchmark_output/tables/action_noncotee_metrics.tex}
\tabcaption{M\'etriques de validation des r\'esidus --- Action non cotee - PE (Black-Scholes apres unsmoothing).}
\subsection{Figures de diagnostic}
\begin{figure}[H]\centering
\includegraphics[width=\linewidth]{benchmark_output/figures/action_noncotee_qq.pdf}
\caption*{\footnotesize\itshape QQ-plot des r\'esidus standardis\'es vs loi normale --- Action non cotee - PE (Black-Scholes apres unsmoothing) (gauche : notre outil ; droite : production).}
\end{figure}
\begin{figure}[H]\centering
\includegraphics[width=\linewidth]{benchmark_output/figures/action_noncotee_acf.pdf}
\caption*{\footnotesize\itshape Autocorr\'elogramme (ACF) des r\'esidus (retards successifs, bande IC 95\,\%) --- Action non cotee - PE (Black-Scholes apres unsmoothing).}
\end{figure}
\clearpage
\section{Immobilier (Black-Scholes apres unsmoothing)}
\subsection{Mod\`ele et donn\'ees}
\textbf{Mod\`ele :} Black-Scholes après unsmoothing AR(1).
\begin{equation*}
\frac{dRE_t}{RE_t}=\mu\,dt+\sigma\,dW_t\quad\text{(série deslissée)}
\end{equation*}
\textbf{Fen\^etre / pas :} 31/12/2005 -- 31/12/2025 \`a pas annuel ($\delta=1$).
\textbf{Donn\'ees :} Indice immobilier, log-rendements annuels deslissés.
\textbf{Note de calibrage :} La dérive $\mu$ production ($\approx5{,}6$) intègre l'hypothèse de réinvestissement des loyers, alors que la série fournie est un indice de prix (rendement en capital seul, dérive $\approx0$). L'écart de $\mu$ est attendu (constat D-1) ; le $R^2$ n'est pas significatif pour un BS.
\subsection{Comparaison des param\`etres}
\input{benchmark_output/tables/immobilier_params.tex}
\tabcaption{Param\`etres calibr\'es : notre outil vs production 2026 --- Immobilier (Black-Scholes apres unsmoothing).}
\subsection{M\'etriques de validation}
\input{benchmark_output/tables/immobilier_metrics.tex}
\tabcaption{M\'etriques de validation des r\'esidus --- Immobilier (Black-Scholes apres unsmoothing).}
\subsection{Figures de diagnostic}
\begin{figure}[H]\centering
\includegraphics[width=\linewidth]{benchmark_output/figures/immobilier_qq.pdf}
\caption*{\footnotesize\itshape QQ-plot des r\'esidus standardis\'es vs loi normale --- Immobilier (Black-Scholes apres unsmoothing) (gauche : notre outil ; droite : production).}
\end{figure}
\begin{figure}[H]\centering
\includegraphics[width=\linewidth]{benchmark_output/figures/immobilier_acf.pdf}
\caption*{\footnotesize\itshape Autocorr\'elogramme (ACF) des r\'esidus (retards successifs, bande IC 95\,\%) --- Immobilier (Black-Scholes apres unsmoothing).}
\end{figure}
\clearpage
\section{Infrastructure (Black-Scholes apres unsmoothing)}
\subsection{Mod\`ele et donn\'ees}
\textbf{Mod\`ele :} Black-Scholes après unsmoothing AR(1).
\begin{equation*}
\frac{dI_t}{I_t}=\mu\,dt+\sigma\,dW_t\quad\text{(série deslissée)}
\end{equation*}
\textbf{Fen\^etre / pas :} 31/12/2005 -- 31/12/2025 \`a pas mensuel ($\delta=1/12$).
\textbf{Donn\'ees :} Indice infrastructure, log-rendements mensuels deslissés.
\textbf{Note de calibrage :} Coefficient d'unsmoothing faible ($b\approx0$) : l'effet du deslissage est marginal. Volatilité du classeur (corrigé de la coquille de facteur 10) en unité moteur (cf. constat U-1).
\subsection{Comparaison des param\`etres}
\input{benchmark_output/tables/infra_params.tex}
\tabcaption{Param\`etres calibr\'es : notre outil vs production 2026 --- Infrastructure (Black-Scholes apres unsmoothing).}
\subsection{M\'etriques de validation}
\input{benchmark_output/tables/infra_metrics.tex}
\tabcaption{M\'etriques de validation des r\'esidus --- Infrastructure (Black-Scholes apres unsmoothing).}
\subsection{Figures de diagnostic}
\begin{figure}[H]\centering
\includegraphics[width=\linewidth]{benchmark_output/figures/infra_qq.pdf}
\caption*{\footnotesize\itshape QQ-plot des r\'esidus standardis\'es vs loi normale --- Infrastructure (Black-Scholes apres unsmoothing) (gauche : notre outil ; droite : production).}
\end{figure}
\begin{figure}[H]\centering
\includegraphics[width=\linewidth]{benchmark_output/figures/infra_acf.pdf}
\caption*{\footnotesize\itshape Autocorr\'elogramme (ACF) des r\'esidus (retards successifs, bande IC 95\,\%) --- Infrastructure (Black-Scholes apres unsmoothing).}
\end{figure}
\clearpage

\section{Conclusion}
La confrontation montre une \textbf{forte coh\'erence} entre notre recalibrage
et les param\`etres en production sur les facteurs \`a retour \`a la moyenne
mod\'elis\'es par maximum de vraisemblance (cr\'edit, dette priv\'ee) : $R^2$
\'elev\'es ($\approx0{,}92$) et param\`etres proches. Les facteurs Vasicek \`a
deux facteurs (inflation, taux r\'eels) pr\'esentent des vitesses de retour et
moyennes long terme coh\'erentes, mais des volatilit\'es diff\'erentes,
imputables \`a la m\'ethodologie de calibrage (cibles distributionnelles en
production vs maximum de vraisemblance), cf.\ constat M-1.

Pour les actions (Hardy) et les actifs r\'eels (Black-Scholes apr\`es
unsmoothing), les moyennes de r\'egime, les matrices de transition \emph{et les
volatilit\'es} se r\'epliquent fid\`element ($\sim$1--2\,\% d'\'ecart) une fois
prise en compte la coquille de facteur 10 (constat U-1), \textbf{depuis
corrig\'ee dans le classeur source}. Le deslissage AR(1) reproduit exactement
la volatilit\'e post-unsmoothing du Private Equity, validant l'impl\'ementation.
L'autocorr\'elation r\'esiduelle d'ordre 1 est, comme attendu, ramen\'ee \`a
z\'ero apr\`es unsmoothing.

Les tests de normalit\'e (Shapiro-Wilk, Jarque-Bera) sont fr\'equemment
rejet\'es, refl\'etant les queues \'epaisses des s\'eries financi\`eres
(\'episodes 2008, 2020, 2022) ; ce constat, commun aux deux jeux de
param\`etres, rel\`eve du choix de mod\`ele (diffusions gaussiennes) et non du
calibrage. Les analyses d\'etaill\'ees et les valeurs num\'eriques figurent
dans les tableaux et figures ci-dessus, reproductibles via le module
\texttt{audit\_mission/benchmark}.

\end{document}
